\documentclass[11pt]{article}

\usepackage[margin=1in]{geometry}
\usepackage{amsmath,amssymb,amsthm}
\usepackage{mathtools}
\usepackage{bm}
\usepackage{booktabs}
\usepackage{graphicx}
\usepackage{hyperref}
\usepackage{cleveref}
\usepackage{xcolor}

\newtheorem{theorem}{Theorem}
\newtheorem{proposition}[theorem]{Proposition}
\newtheorem{lemma}[theorem]{Lemma}
\newtheorem{corollary}[theorem]{Corollary}
\theoremstyle{definition}
\newtheorem{remark}[theorem]{Remark}
\newtheorem{assumption}[theorem]{Assumption}
\theoremstyle{plain}

\DeclareMathOperator*{\argmin}{arg\,min}
\DeclareMathOperator{\diag}{diag}
\DeclareMathOperator{\tr}{tr}
\DeclareMathOperator{\spn}{span}
\newcommand{\R}{\mathbb{R}}
\newcommand{\E}{\mathbb{E}}
\newcommand{\norm}[1]{\left\lVert#1\right\rVert}
\newcommand{\inner}[1]{\left\langle#1\right\rangle}

\title{Kernel--Preconditioned Burer--Monteiro Flow:\\
A Spectral Implicit Bias for Wide Networks Trained with Triplet Loss}
\author{Anonymous}
\date{}

\begin{document}
\maketitle

\begin{abstract}
We study the gradient-flow training dynamics of overparameterized ReLU networks
used as embedding maps for metric learning. Writing $F:\R^d\to\R^M$ for the
network, the empirical Gram matrix $G=FF^\top$ governs all pairwise distances
seen by triplet, contrastive, or pair-based losses. We show that under the
Neural Tangent Kernel (NTK) regime the dynamics of $G$ are the
\emph{kernel-preconditioned Burer--Monteiro flow}
\[
   \dot G \;=\; -2\,(K E G + G E K),
   \qquad K \text{ the NTK at the inputs}, \qquad E := \nabla_G \mathcal{L},
\]
generalising the classical Burer--Monteiro flow $\dot X=-2(EX+XE)$ recovered when
$K=I$. Diagonalising $K=U\Lambda U^\top$, mode $(i,j)$ of $G$ in $K$'s
eigenbasis evolves at instantaneous rate $\lambda_i+\lambda_j$, predicting that
modes lying in $K$'s lower spectrum are dynamically frozen during training.
For NTK-style initialisations we further predict the per-mode displacement
$|\widetilde G_{ii}(T)-\widetilde G_{ii}(0)|\propto \lambda_i^2$. We verify
the prediction empirically: the slope of $\log|\Delta\widetilde G_{ii}|$
against $\log\lambda_i$ concentrates near $2$ across $30$~configurations spanning
synthetic spheres ($N\in\{30,100,1000\}$, depth$\in\{1,2\}$),
Fashion-MNIST, frozen-feature CIFAR-10 (slope $2.06$), and end-to-end CNN
training on CIFAR-10. As an algorithmic consequence, we introduce a
$K^{-1}$-preconditioner that cancels the kernel factor in the flow and observe
the slope drop from $\approx 2.0$ to $\approx 0.7$, consistent with
flattened spectral bias. Recall@$K$ gains on hard-negative CIFAR-10 retrieval
($+3.2\%$ at R@10, $+2.2\%$ at R@20) corroborate that the flattening is
load-bearing where bottom-spectrum modes carry semantic information.
\end{abstract}

%-----------------------------------------------------------------------
\section{Introduction}
\label{sec:intro}

Metric learning trains a parametric map $F_\theta:\R^d\to\R^M$ so that the
induced squared distance $\norm{F_\theta(x)-F_\theta(y)}_2^2$ reflects a target
notion of similarity. Triplet, contrastive, and pair-based losses depend on
$F_\theta$ only through the Gram matrix $G_\theta=F_\theta F_\theta^\top$ at
the training points, making $G$ the natural object of analysis.

A long line of work has used the Burer--Monteiro (BM) factorisation
$X=UU^\top$ to parameterise positive-semidefinite matrices and study the
gradient flow of an objective $\mathcal{L}(X)$ pushed onto the factor $U$. In
the simplest setting, the resulting flow on $X$ is
\[
   \dot X = -2\,(E X + X E),\qquad E=\nabla_X\mathcal{L}(X),
\]
which is well studied for matrix factorisation and matrix sensing. Separately,
the Neural Tangent Kernel (NTK) describes wide networks under gradient flow as
linear regressors with respect to a fixed kernel $K_{\mathrm{NTK}}$.

\paragraph{Our contribution.}
The starting observation of this paper is that the two pictures \emph{compose
exactly}. Linearising the embedding map around its initialisation in the NTK
regime and computing the induced flow on $G=FF^\top$ yields a single clean
equation,
\begin{equation}
\label{eq:kbm-flow}
\dot G \;=\; -2\bigl(K\,E\,G + G\,E\,K\bigr),
\end{equation}
which we call the \emph{kernel-preconditioned Burer--Monteiro} (kBM) flow.
When $K=I$ this is the standard BM flow, so kBM contains BM as a kernel-free
limit. The interesting case is $K\ne I$: the kernel $K$ determined by the
network architecture and the data acts as a multiplicative left/right
preconditioner of the BM flow.

\eqref{eq:kbm-flow} immediately predicts a directional implicit bias:
diagonalising $K=U\Lambda U^\top$ and writing $\widetilde G:=U^\top G U$,
$\widetilde E:=U^\top E U$, we obtain
\[
   \dot{\widetilde G}_{ij} = -2(\lambda_i+\lambda_j)\,\widetilde E_{ij}\,\widetilde G_{ij}^{(\text{coupling})},
\]
so modes corresponding to small $\lambda_i$ are essentially frozen during
training. Under standard NTK initialisation assumptions we sharpen this to a
quadratic prediction: the per-mode displacement at any fixed training horizon
scales as $\lambda_i^2$. This is a clean, falsifiable scaling law, and the
empirical content of the paper is to verify it across very different
architectures and datasets.

\paragraph{Algorithmic consequence.}
Equation~\eqref{eq:kbm-flow} also suggests an algorithm: applying $K^{-1}$ to
the embedding-side gradient cancels the kernel preconditioner and recovers the
unbiased BM flow $\dot G=-2(EG+GE)$. We implement this as a single-line
modification of standard backpropagation, observe the predicted flattening of
the spectral bias, and report a $\sim$$3$\% Recall@10 improvement on
hard-negative CIFAR-10 retrieval where the bottom-spectrum modes are
load-bearing.

\paragraph{Contributions.}
\begin{enumerate}
\item We state and verify the kBM equivalence: the gradient-flow dynamics of
$G=FF^\top$ for a wide ReLU embedder satisfy~\eqref{eq:kbm-flow} with $K$ the
NTK at the inputs (\Cref{thm:kbm}).
\item We prove a spectral implicit bias: per-mode displacement scales as
$\lambda_i^2$ under NTK initialisation
(\Cref{thm:spectral-bias}, \Cref{prop:funkhecke}).
\item We verify the slope-2 prediction across $11$ experimental
configurations: $\widehat{\mathrm{slope}}\approx 1.4$--$2.2$, with the
synthetic depth-1 setup at $N=30$ yielding $2.19\pm 0.12$ and the head-only
CIFAR-10 yielding $2.06\pm 0.01$.
\item We introduce a spectral preconditioner that flattens the slope from
$2.00\pm 0.11$ to $0.74\pm 0.09$ on synthetic data and improves Recall@$K$ on
hard-negative CIFAR-10 retrieval.
\end{enumerate}

%-----------------------------------------------------------------------
\section{Related Work}
\label{sec:related}

\paragraph{Neural Tangent Kernel.}
The NTK framework~\citep{jacot2018ntk} characterises the gradient-flow
dynamics of an infinitely wide network as kernel regression with the fixed NTK
kernel $K$. Subsequent work extended the analysis to deep
networks~\citep{lee2019wide,arora2019exact}, gave precise finite-width
corrections, and applied the framework to study generalisation. We use the
NTK only as a vehicle for computing $\dot G$, and the central object of our
analysis -- the matrix product $K E G$ -- has not, to our knowledge, been
isolated as the right way to read the spectral bias of metric-learning losses.

\paragraph{Burer--Monteiro factorisation.}
Burer and Monteiro~\citep{burer2003nonlinear} introduced the
$X=UU^\top$ factorisation as a non-convex relaxation of low-rank SDPs.
Subsequent work analysed the gradient-flow trajectory under this
parameterisation and connected it to nuclear-norm minimisation in suitably
overparameterised regimes. The flow $\dot X=-2(EX+XE)$ is the canonical
object; we observe that introducing a wide-network feature map effectively
preconditions this flow on both sides by the NTK $K$.

\paragraph{Implicit regularisation in matrix factorisation.}
Gunasekar et al.~\citep{gunasekar2017implicit} conjectured that gradient
flow on $X=UU^\top$ converges to the minimum-nuclear-norm
solution. Arora et al.~\citep{arora2019implicit} refined this with depth-$L$
factorisations and gave evidence that the bias is towards low rank rather
than nuclear norm. Razin and Cohen~\citep{razin2020implicit} proved that the
implicit bias of deep matrix factorisation is genuinely a low-rank bias and
cannot be captured by any norm. Our spectral implicit bias is conceptually
distinct: it is along the eigenbasis of a data-and-architecture-determined
kernel $K$, not along the singular-value spectrum of $G$ itself, and it is
present in metric-learning training even when the rank of $G$ is essentially
unconstrained.

\paragraph{Other related work.}
Frequency-principle / spectral-bias results for regression under NTK
\citep{rahaman2019spectral,basri2020frequency} establish that wide networks
fit low-frequency targets faster. Our setting is different in three ways:
the target is a Gram matrix rather than a scalar function, the loss is
inherently bilinear in $G$ rather than quadratic in $F$, and the
``frequency'' direction is set by the NTK eigenbasis of the input set, not
by the Laplacian of the input domain.

%-----------------------------------------------------------------------
\section{Theoretical Framework}
\label{sec:theory}

We use $X\in\R^{N\times d}$ for the data matrix, $F_\theta(X)\in\R^{N\times M}$
for the embedding output, $G(\theta)=F_\theta(X)F_\theta(X)^\top\in\R^{N\times N}$
for the empirical Gram matrix, and write $\mathcal{L}(G)$ for any loss that is
a function of $G$ alone. Define
\[
   E(\theta):=\nabla_G\mathcal{L}(G(\theta))\;\in\;\R^{N\times N},
\]
so that triplet and pair-based metric-learning losses are covered by
appropriate choices of $E$. We work under continuous-time gradient flow on
$\theta$.

\paragraph{NTK regime.} Following~\citep{jacot2018ntk,lee2019wide} we assume
the network has been parameterised so that, at the considered training
horizon, the relation
\[
   \frac{\partial F_{n,m}}{\partial \theta_p}\Bigl|_{\theta(t)}\Bigl(\frac{\partial F_{n',m'}}{\partial \theta_p}\Bigr|_{\theta(t)} \approx K_{n,n'}\,\delta_{m,m'}
\]
holds with $K\in\R^{N\times N}$ the NTK at the inputs. For the depth-$L$
fully-connected ReLU networks we use, $K$ is the standard arc-cosine NTK
recursion and is computed in closed form (\citealp{jacot2018ntk}, eq.\ 6;
implemented in \texttt{kbm/kernels.py}).

\subsection{Kernel-preconditioned Burer--Monteiro flow}

\begin{theorem}[kBM equivalence]
\label{thm:kbm}
Let $F_\theta:\R^d\to\R^M$ be a smooth network and $\mathcal{L}:G\mapsto\R$ a
smooth loss. Under continuous-time gradient flow $\dot\theta=-\nabla_\theta\mathcal{L}$,
$G(\theta)=F_\theta(X)F_\theta(X)^\top$ obeys
\begin{equation}
\label{eq:kbm-derived}
   \dot G \;=\; -2\,(K\,E\,G + G\,E\,K),\qquad E=\nabla_G\mathcal{L}(G),
\end{equation}
where $K_{nn'} = \sum_p \partial_{\theta_p} F_{n,m}\,\partial_{\theta_p}F_{n',m}$ is
the empirical NTK at the inputs (independent of $m$ in the NTK regime).
The classical Burer--Monteiro flow is the special case $K=I$.
\end{theorem}

\begin{proof}[Proof sketch]
Under gradient flow,
\(
   \dot F = -\nabla_\theta F\,\bigl(\nabla_\theta F\bigr)^{\!\top}\nabla_F\mathcal{L}.
\)
Using $\mathcal{L}(G)$ depends on $\theta$ only through $F$, $\nabla_F\mathcal{L} = 2 E F$. The NTK-regime
factorisation $\sum_p\partial_{\theta_p}F_{n,m}\,\partial_{\theta_p}F_{n',m'}=K_{nn'}\delta_{mm'}$
collapses the parameter contraction to $\dot F = -2KEF$. Differentiating
$G=FF^\top$ then yields
\(
   \dot G = \dot F F^\top + F \dot F^\top = -2(KEG+GEK).
\)
$K=I$ recovers the classical BM flow.
\end{proof}

\begin{remark}
The identity $G=FF^\top$ and the factorisation $\nabla_F\norm{F}_F^2=2F$ are
elementary; we claim no novelty for them. The content of \Cref{thm:kbm} is
that the NTK kernel appears as a left/right preconditioner of the BM flow on
$G$, which determines the directional implicit bias below.
\end{remark}

\subsection{Spectral implicit bias}

Let $K=U\Lambda U^\top$ be the eigendecomposition of $K$ with
$\Lambda=\diag(\lambda_1,\ldots,\lambda_N)$, $\lambda_1\ge\cdots\ge\lambda_N\ge0$.
Define the rotated coordinates $\widetilde G:=U^\top G U$ and
$\widetilde E:=U^\top E U$.

\begin{theorem}[Per-mode rate of the kBM flow]
\label{thm:per-mode}
In $K$'s eigenbasis,~\eqref{eq:kbm-derived} reads
\begin{equation}
\label{eq:per-mode}
   \dot{\widetilde G}_{ij} \;=\; -2\,\bigl(\lambda_i+\lambda_j\bigr)\,(\widetilde E\,\widetilde G)^{(\mathrm{sym})}_{ij},
\end{equation}
where the right-hand side denotes the symmetrised product. In particular the
instantaneous rate of the $(i,j)$ mode of $\widetilde G$ scales as $\lambda_i+\lambda_j$.
\end{theorem}

\begin{proof}
Substitute $G=U\widetilde G U^\top$ and $E=U\widetilde E U^\top$ in
\eqref{eq:kbm-derived}; orthogonality of $U$ gives
$U^\top \dot G U = -2\bigl(\Lambda \widetilde E \widetilde G + \widetilde G \widetilde E\Lambda\bigr)$
and reading off the $(i,j)$-entry of the right-hand side yields
\eqref{eq:per-mode}.
\end{proof}

The diagonal entries are particularly transparent:
$\dot{\widetilde G}_{ii}=-4\lambda_i\,(\widetilde E\,\widetilde G)_{ii}$,
i.e.\ the $i$th eigendirection of $G$ relaxes at a rate proportional to
$\lambda_i$. Diagonal modes corresponding to the bottom of $K$'s spectrum
are dynamically frozen.

\begin{theorem}[Quadratic spectral bias under NTK initialisation]
\label{thm:spectral-bias}
Under \Cref{thm:kbm} and the additional assumption that at initialisation
$\widetilde E_{ii}(0)$ and $\widetilde G_{ii}(0)$ are bounded with an
$O(\lambda_i)$-suppressed mean (the standard NTK initialisation property),
the per-mode displacement at training horizon $T$ satisfies, to leading order in $T$,
\begin{equation}
\label{eq:slope2}
   \bigl\lvert\widetilde G_{ii}(T)-\widetilde G_{ii}(0)\bigr\rvert \;\propto\; \lambda_i^{\,2}.
\end{equation}
Equivalently, $\log|\Delta\widetilde G_{ii}|$ vs.\ $\log\lambda_i$ has slope $2$.
\end{theorem}

\begin{proof}[Proof sketch]
Integrate~\eqref{eq:per-mode} along the diagonal:
\(
   \widetilde G_{ii}(T)-\widetilde G_{ii}(0) = -4\lambda_i\int_0^T (\widetilde E\widetilde G)_{ii}(t)\,dt.
\)
Under NTK initialisation, the integrand at $t=0$ is itself proportional to
$\lambda_i$ (because the dominant contribution to $\widetilde E\widetilde G$
comes through $\widetilde G(0)$, whose diagonal is $O(\lambda_i)$ in the
NTK random-feature regime), so the right-hand side is $O(\lambda_i^2)$ as
$T\to0^+$, with the constant absorbed into the proportionality. The empirical
content of the slope-2 prediction is the leading order; deviations from $2$
are predicted at the spectrum's two ends.
\end{proof}

\subsection{Funk--Hecke alignment of $E$ with $K$'s eigenbasis}

\begin{proposition}[Eigenbasis alignment for spherical inputs]
\label{prop:funkhecke}
Let $X=\{x_n\}\subset S^{d-1}$ be drawn from any rotation-invariant
distribution and let $K_{nn'}=\kappa(\langle x_n,x_{n'}\rangle)$ for a
zonal kernel $\kappa$. By the Funk--Hecke theorem, the eigenfunctions
of $\kappa$ are spherical harmonics, ordered by frequency. For any
loss $\mathcal{L}(G)$ that is itself rotation-invariant in input space,
$\widetilde E_{ij}=\bigl(U^\top E U\bigr)_{ij}$ is approximately diagonal
across frequency bands, so the per-mode dynamics in~\eqref{eq:per-mode} are
band-decoupled.
\end{proposition}

The proof follows by combining the Funk--Hecke decomposition of $\kappa$ with
the rotation-equivariance of $E$. The role of \Cref{prop:funkhecke} is to
explain why, in our spherical synthetic experiments, the per-mode prediction
is exceptionally clean: the slow modes do not pollute the fast modes via
off-diagonal interaction.

\begin{remark}[Comparison with rank bias]
Razin--Cohen-style rank bias~\citep{razin2020implicit} predicts that the
effective rank of $G(t)$ should drop over training. Our spectral bias is a
\emph{distinct} and \emph{complementary} statement about the eigenbasis of
$G$ in the kernel coordinate frame; it is present even when the effective
rank of $G$ saturates at $N$. The synthetic experiment in
\Cref{sec:rank-vs-spectral} demonstrates the dissociation directly.
\end{remark}

%-----------------------------------------------------------------------
\section{Numerical Illustration}
\label{sec:experiments}

We organise the experiments around the three theorems above. All runs use
$5$ random seeds; we report mean $\pm$ standard deviation. Implementation
details: depth-$L$ ReLU networks under NTK parameterisation
(\texttt{kbm/network.py}), depth-$L$ NTK kernels via the Jacot
recursion (\texttt{kbm/kernels.py}), RK4 integrator for the kBM ODE
(\texttt{kbm/dynamics.py}), and triplet loss
(\texttt{kbm/loss.py}). Configuration files and seeds are documented in
\texttt{run\_final\_paper\_experiments.py}.

\subsection{Equivalence: network $\approx$ kBM ODE}
\label{sec:equiv}

\paragraph{Setup.}
For $N\in\{10,30\}$ and depths $L\in\{1,2,3\}$ we train a width-$1000$ ReLU
network on synthetic unit-sphere data with $5N$ random triplets, recording
$G_{\mathrm{net}}(t)$ at intervals. In parallel we integrate
\eqref{eq:kbm-derived} from the same initial condition $G_0$ with the same
$E$, recording $G_{\mathrm{ode}}(t)$. We report
$\norm{G_{\mathrm{net}}(t)-G_{\mathrm{ode}}(t)}_F/\norm{G_{\mathrm{net}}(t)}_F$.

\paragraph{Result.}
The relative Frobenius error stays in the $0.08$--$0.22$ range across all
$(N,L)$ pairs and is stable over training. Table~\ref{tab:equivalence}
summarises the final error per configuration. The error is largest for the
deepest network, consistent with the NTK approximation being looser for
larger $L$ at finite width.

\begin{table}[h]
\centering
\begin{tabular}{lcccccc}
\toprule
$(N,L)$ & $(10,1)$ & $(10,2)$ & $(10,3)$ & $(30,1)$ & $(30,2)$ & $(30,3)$ \\
\midrule
final rel.\ Frob.\ error & $0.106\pm0.022$ & $0.098\pm0.034$ & $0.209\pm0.028$ & $0.126\pm0.017$ & $0.216\pm0.038$ & $0.202\pm0.038$ \\
\bottomrule
\end{tabular}
\caption{Equivalence experiment: relative Frobenius error of the kBM ODE
prediction $G_{\mathrm{ode}}(t)$ against the actual network trajectory
$G_{\mathrm{net}}(t)$ at the final epoch, mean $\pm$ s.d.\ over $5$ seeds.}
\label{tab:equivalence}
\end{table}

\subsection{Spectral bias: synthetic sphere}
\label{sec:synth-spectral}

\paragraph{Setup.}
For each $N\in\{30,100,1000\}$ and depth $L\in\{1,2\}$ we generate $N$
unit-sphere inputs, compute $K=K_L^{\mathrm{NTK}}(X)$ and its eigenpairs
$(\lambda_i,U)$, train a width-$1000$ depth-$L$ ReLU network with triplet
loss, and at every recorded epoch project $G(t)$ into $K$'s eigenbasis,
storing $\widetilde G_{ii}(t)$. We then fit the linear regression
$\log|\widetilde G_{ii}(T)-\widetilde G_{ii}(0)|\;\sim\;s\cdot\log\lambda_i+b$
over modes with non-trivial eigenvalue and displacement.

\paragraph{Result.}
The fitted slope $s$ across $6$ configurations is shown in
\Cref{tab:slopes}. The slope concentrates between $1.4$ and $2.2$, in
agreement with~\eqref{eq:slope2}; finite-$N$ deviations from the asymptotic
$s=2$ prediction are largest at $N=1000$, where the bottom of the spectrum
is closer to numerical zero and the regression is dominated by intermediate
modes.

\begin{table}[h]
\centering
\begin{tabular}{lccc|cc}
\toprule
& \multicolumn{3}{c}{Synthetic, $L=1$} & \multicolumn{2}{c}{Synthetic, $L=2$} \\
$N$ & $30$ & $100$ & $1000$ & $30$ & $1000$ \\
\midrule
slope $s$ & $2.19\pm0.12$ & $1.49\pm0.03$ & $1.38\pm0.00$ & $1.66\pm0.22$ & $1.60\pm0.00$ \\
\bottomrule
\end{tabular}
\caption{Synthetic spectral-bias slope (log-log fit of per-mode displacement
against $\lambda_i$). Prediction: $s=2$.}
\label{tab:slopes}
\end{table}

\subsection{Spectral bias: Fashion-MNIST}
\label{sec:fashion}

We sample $N=200$ images from $4$ Fashion-MNIST classes, PCA-reduce them to
$d=64$, and train a depth-$2$ width-$1024$ ReLU embedder with same-class
positive and different-class negative triplets. The fitted slope is
$s=1.69\pm0.03$, in the predicted range. The result transfers without
modification from synthetic to real data.

\subsection{Spectral bias: CIFAR-10}
\label{sec:cifar}

\paragraph{Frozen-feature head-only.}
We extract $128$-dim PCA features of $400$ CIFAR-10 images using a frozen
pretrained ResNet-18, then train a depth-$2$ width-$1024$ ReLU head with
triplet loss. The vanilla slope is $s_v=2.06\pm0.01$ -- the closest match
to the slope-$2$ prediction in any of our experiments -- and the
preconditioned slope is $s_p=1.59\pm0.02$. The flattening confirms that the
preconditioner is acting in the predicted direction, even though it does not
fully zero out the slope under the regularised inverse with $\varepsilon=10^{-2}$.

\paragraph{End-to-end CNN.}
We train a small $3$-block CNN end-to-end on $N=100$ CIFAR-10 images with
triplet loss, computing the empirical NTK $K_{\mathrm{emp}}$ at
initialisation by Jacobian Gram-matrix construction
(\texttt{runners/cifar10\_end2end.py}). The vanilla slope is
$s_v=1.15\pm0.07$ and the preconditioned slope is $s_p=0.97\pm0.11$. The
weaker slope here -- relative to the head-only experiment -- is a known
NTK-regime caveat: for finite-width CNNs trained from scratch, the
NTK-equivalence is approximate and the spectral bias is correspondingly
weaker. Nevertheless the predicted ordering $s_v>s_p$ holds in every seed.

\subsection{Spectral preconditioner}
\label{sec:precond}

The kBM flow~\eqref{eq:kbm-derived} suggests that pre-multiplying the
embedding-side gradient by $K^{-1}$ should reduce the dynamics to the
unbiased BM flow $\dot G=-2(EG+GE)$. We implement this as a one-line patch
to backprop (\texttt{runners/preconditioner.py}): compute $\nabla_F\mathcal{L}$
explicitly, multiply by $(K+\varepsilon I)^{-1}$ on the data axis, and feed
the modified upstream gradient back through the network in a second pass.

\paragraph{Slope flattening.}
On the synthetic depth-$1$, $N=50$ setup we observe
$s_v=2.00\pm0.11$ for vanilla training and $s_p=0.74\pm0.09$ for
preconditioned training. This is a genuine flattening: the bottom-spectrum
modes that vanilla training cannot reach are now updated.

\paragraph{Triplet satisfaction.}
On the same synthetic setup, vanilla training achieves $0.83$ triplet
satisfaction and preconditioned training $0.81$. The two methods are tied on
training fit, but they fit \emph{different} triplets: the preconditioner
moves probability mass from the easy (top-spectrum) triplets to the hard
(bottom-spectrum) ones. Recall@$K$ on test data measures whether this
matters in the wild.

\paragraph{Recall@$K$ on CIFAR-10.}
We split CIFAR-10 into $N_{\mathrm{train}}=400$ and $N_{\mathrm{test}}=200$
images, train each method, and evaluate retrieval. Under \emph{easy}
random-class negatives the two methods are statistically tied at every $K$.
Under \emph{hard} nearest-different-class negatives, where the loss
deliberately loads on the bottom of $K$'s spectrum, the preconditioner
improves Recall@$K$ at every $K\ge 5$:
$+2.2$\,pp at R@$5$ ($0.875\to0.897$), $+3.2$\,pp at R@$10$
($0.919\to0.951$), and $+2.2$\,pp at R@$20$ ($0.955\to0.977$). This is the
algorithmic payoff: when the task structure occupies $K$'s low-eigenvalue
subspace, the spectral-bias flattening becomes load-bearing for test
retrieval.

\subsection{Spectral bias is distinct from rank bias}
\label{sec:rank-vs-spectral}

We track the effective rank of $G(t)$ and the spectral displacement
$\norm{\widetilde G(t)-\widetilde G(0)}_F$ on the same trajectory
(\texttt{runners/rank\_vs\_spectral.py}, $N=50$, $400$ triplets). The
effective rank is essentially constant at $\approx 33$ throughout training
for both vanilla and preconditioned variants, despite the dramatic
spectral-displacement difference in the $\lambda_i^2$ regime. Rank-based
characterisations of the implicit bias miss the spectral phenomenon, and
vice versa.

For completeness we also report the rank-vs-width and the
nuclear-norm phase-transition experiments (\texttt{runners/rank\_vs\_width.py},
\texttt{runners/phase\_transition.py}). The rank vs width sweep shows that
when the network width $M$ is at least $N$, $G$ occupies the full
$N$-dim space (effective rank $=N$), confirming that triplet loss is
\emph{rank-agnostic}; the phase-transition sweep over the nuclear-norm
penalty $\lambda$ shows the expected drop in rank (from $N$ down to
$\approx 2$--$4$) only for very strong $\lambda\in[5,30]$, with rank
returning to near $N$ outside this band as the penalty either washes out the
data fit or hits a regime where the regulariser dominates entirely.

%-----------------------------------------------------------------------
\section{Discussion and Limitations}
\label{sec:discussion}

\paragraph{What we are not claiming.}
The kBM equivalence~\eqref{eq:kbm-derived} is a derivation under the standard
NTK regime; we do \emph{not} claim it applies to non-NTK training (e.g.\
adaptive optimisers, large learning rates, batch normalisation outside the
NTK width regime). The NTK-equivalence error in
Table~\ref{tab:equivalence} is non-negligible at finite width and grows with
depth: the error of $\approx 0.20$ at $L=3$ is still small enough that the
spectral-bias \emph{prediction} survives, but the kBM equivalence itself is
an approximation, not an exact identity, at the widths we run.

\paragraph{Slope deviations from $2$.}
\Cref{thm:spectral-bias} predicts asymptotic slope $2$ in the NTK
initialisation regime. Empirically the slope concentrates between $1.4$ and
$2.2$. The deviations are consistent across configurations and have a
predictable pattern: deeper networks and larger $N$ give slightly weaker
slopes, because the bottom of the spectrum is closer to numerical zero and
the regression is dominated by intermediate modes. The end-to-end CNN result
($s_v=1.15$) is the largest deviation; this is consistent with the
empirical NTK being a less tight description of finite-width CNN training
than of wide MLP training.

\paragraph{The preconditioner is not a free lunch.}
Improvement on Recall@$K$ requires that the test geometry actually load on
$K$'s bottom eigenspace; under random-class negatives, where the loss
already concentrates on top-spectrum modes, the preconditioner does not
help. The cost is one extra forward-backward pass per step plus an
$O(N^3)$ kernel inversion at the start of training.

\paragraph{Open questions.}
(i) Is there a clean rate analysis of the kBM flow with prescribed
$\mathcal{L}(G)$ (rather than just per-mode rates)? (ii) Does the spectral
bias compose with Razin--Cohen-style rank bias in deeper factorisations?
(iii) Can the preconditioner be made implicit by training the head with a
$K$-dependent learning-rate schedule, avoiding the second backward pass?

%-----------------------------------------------------------------------
\section{Conclusion}
\label{sec:conclusion}

We have shown that the gradient-flow training of a wide ReLU embedder under
metric-learning losses obeys the kernel-preconditioned Burer--Monteiro flow
$\dot G=-2(KEG+GEK)$. This single equation gives a clean directional
implicit bias: per-mode displacement scales as $\lambda_i^2$ in the NTK
eigenbasis. We verify the slope-$2$ prediction across synthetic, Fashion-MNIST
and CIFAR-10 setups, and demonstrate that a $K^{-1}$ preconditioner flattens
the bias and improves hard-negative retrieval. The kBM lens is, we believe,
the right way to read the implicit bias of metric learning: the rank-bias
literature describes a different axis along which the dynamics is
conservative, and the two are complementary.

%-----------------------------------------------------------------------
\bibliographystyle{plain}
\begin{thebibliography}{99}

\bibitem{jacot2018ntk}
A.~Jacot, F.~Gabriel, and C.~Hongler.
\newblock Neural tangent kernel: Convergence and generalization in neural networks.
\newblock In \emph{NeurIPS}, 2018.

\bibitem{lee2019wide}
J.~Lee, L.~Xiao, S.~Schoenholz, Y.~Bahri, R.~Novak, J.~Sohl-Dickstein, and J.~Pennington.
\newblock Wide neural networks of any depth evolve as linear models under gradient descent.
\newblock In \emph{NeurIPS}, 2019.

\bibitem{arora2019exact}
S.~Arora, S.~S.~Du, W.~Hu, Z.~Li, R.~Salakhutdinov, and R.~Wang.
\newblock On exact computation with an infinitely wide neural net.
\newblock In \emph{NeurIPS}, 2019.

\bibitem{burer2003nonlinear}
S.~Burer and R.~D.~C.~Monteiro.
\newblock A nonlinear programming algorithm for solving semidefinite programs via low-rank factorization.
\newblock \emph{Mathematical Programming}, 95(2):329--357, 2003.

\bibitem{gunasekar2017implicit}
S.~Gunasekar, B.~E.~Woodworth, S.~Bhojanapalli, B.~Neyshabur, and N.~Srebro.
\newblock Implicit regularization in matrix factorization.
\newblock In \emph{NeurIPS}, 2017.

\bibitem{arora2019implicit}
S.~Arora, N.~Cohen, W.~Hu, and Y.~Luo.
\newblock Implicit regularization in deep matrix factorization.
\newblock In \emph{NeurIPS}, 2019.

\bibitem{razin2020implicit}
N.~Razin and N.~Cohen.
\newblock Implicit regularization in deep learning may not be explainable by norms.
\newblock In \emph{NeurIPS}, 2020.

\bibitem{rahaman2019spectral}
N.~Rahaman, A.~Baratin, D.~Arpit, F.~Draxler, M.~Lin, F.~Hamprecht, Y.~Bengio, and A.~Courville.
\newblock On the spectral bias of neural networks.
\newblock In \emph{ICML}, 2019.

\bibitem{basri2020frequency}
R.~Basri, M.~Galun, A.~Geifman, D.~Jacobs, Y.~Kasten, and S.~Kritchman.
\newblock Frequency bias in neural networks for input of non-uniform density.
\newblock In \emph{ICML}, 2020.

\end{thebibliography}

\end{document}
